\begin{abstract}
Removing a circuit motif need not remove the computation it supports. On the
Drosophila MaleCNS connectome, two reconstructed graph interventions preserve
identical degree and module-pair statistics while eliminating all designated
common-relay triples. Yet one retains both task sources and the other loses
access to one, consistent with archived classification accuracies of 97.45\%
and 50.55\%. Readout retraining on the reconstructed transfers reproduces this
high-versus-chance contrast, without recovering the identical historical run
instances. To explain what survives a graph intervention, we introduce
task-information auditing, which separates structural constraints, information
available at the readout, and finite-model behavior through a compact
source-to-readout transfer matrix. For a fixed two-step nonnegative linear
propagator, we prove that perfect Bayes solvability of a noiseless subspace
XOR task requires neither a common relay nor full transfer rank, but exactly
the visibility of both sources. Across two additional graph extractions, ideal
solvability remains perfect at every nonzero intervention dose, even as
trained accuracy rises from approximately 67\% to 94\%. In a four-source
CIFAR-10 probe, an independent transfer direction accompanies an increase in
target-source classification accuracy from 16.4\% to 56.6\%. These findings
shift the evaluation of graph interventions from the structures they preserve
to the task-relevant computations they leave intact.
\end{abstract}
